\section{Experiments}
\label{sec:experiments}

\subsection{Experimental Setup}
\label{sec:setup}

\paragraph{Baselines.}
We evaluate seven continual learning methods spanning the spectrum from no forgetting mitigation to multimodal-aware strategies:
\begin{itemize}
    \item \textbf{Naive Sequential:} Fine-tunes on each task sequentially without any forgetting mitigation (lower bound).
    \item \textbf{Joint Training:} Trains on the union of all tasks simultaneously (upper bound for knowledge retention, though not for new knowledge acquisition).
    \item \textbf{EWC}~\citep{Kirkpatrick2017}: Elastic Weight Consolidation penalizes changes to parameters important for previous tasks via the Fisher information matrix ($\lambda = 10$).
    \item \textbf{Experience Replay}~\citep{Rolnick2019}: Maintains a memory buffer of exemplars from previous tasks and replays them during new task training (buffer size = 10\% of each task).
    \item \textbf{LKGE}~\citep{Cui2023}: A lifelong KG embedding framework with embedding regularization and knowledge distillation, using TransE as the base model.
    \item \textbf{RAG Agent}~\citep{Lewis2020}: A retrieval-augmented generation pipeline that retrieves relevant KG subgraphs and uses Qwen2.5-7B as the language model for KGQA.
    \item \textbf{CMKL (ours):} Cross-Modal Knowledge Learner with modality-aware EWC, multimodal memory replay, and cross-attention fusion using DistMult as the decoder.
\end{itemize}

\paragraph{Hyperparameters.}
All KGE-based methods use embedding dimension $d = 256$, learning rate $\eta = 0.001$, batch size 512, and are trained for 100 epochs per task.
We use the Adam optimizer.
For EWC, $\lambda = 10$; for experience replay, the buffer stores 500 exemplars per task with relation-balanced selection.
CMKL uses a 4-head cross-attention fusion module, modality-specific Fisher matrices, and K-means diverse selection for the multimodal memory buffer (1000 exemplars).

\paragraph{Infrastructure.}
All experiments are conducted on NVIDIA V100 GPUs (32\,GB) on the KAUST IBEX cluster.
Each method is trained with five random seeds $\{42, 123, 456, 789, 1024\}$ and we report mean $\pm$ standard deviation.
Statistical significance is assessed via paired $t$-tests at $p < 0.05$.

\subsection{Link Prediction Results}
\label{sec:lp_results}

\Cref{tab:main_results} presents the main link prediction results across all seven methods.

\begin{table}[t]
\centering
\caption{Link prediction results on \ours (mean $\pm$ std over 5 seeds). AP: Average Performance (higher is better), AF: Average Forgetting (lower is better), BWT: Backward Transfer (higher is better), REM: Remembering (higher is better). Best AP is \textbf{bolded}.}
\label{tab:main_results}
\small
\begin{tabular}{lcccc}
\toprule
\textbf{Method} & \textbf{AP} $\uparrow$ & \textbf{AF} $\downarrow$ & \textbf{BWT} $\uparrow$ & \textbf{REM} $\uparrow$ \\
\midrule
Naive Sequential & $0.004 \pm 0.000$ & $0.021 \pm 0.000$ & $-0.021 \pm 0.000$ & $0.980 \pm 0.000$ \\
Joint Training & $0.018 \pm 0.000$ & $0.000 \pm 0.000$ & $0.000 \pm 0.000$ & $1.000 \pm 0.000$ \\
EWC~\citep{Kirkpatrick2017} & $0.004 \pm 0.000$ & $0.017 \pm 0.001$ & $-0.017 \pm 0.001$ & $0.984 \pm 0.001$ \\
Experience Replay~\citep{Rolnick2019} & $0.004 \pm 0.000$ & $0.021 \pm 0.000$ & $-0.021 \pm 0.000$ & $0.980 \pm 0.000$ \\
LKGE~\citep{Cui2023} & $0.039 \pm 0.001$ & $0.012 \pm 0.002$ & $-0.010 \pm 0.003$ & $0.990 \pm 0.003$ \\
RAG (Qwen2.5-7B) & $0.002 \pm 0.001$ & $0.001 \pm 0.001$ & $-0.001 \pm 0.001$ & $0.999 \pm 0.001$ \\
\midrule
\textbf{CMKL (ours)} & $\textbf{0.063} \pm 0.003$ & $0.040 \pm 0.003$ & $-0.040 \pm 0.003$ & $0.960 \pm 0.003$ \\
\bottomrule
\end{tabular}
\end{table}


\paragraph{CMKL achieves the highest average performance.}
CMKL with the DistMult decoder achieves AP\,=\,0.063$\pm$0.003, substantially outperforming all other methods.
This is 3.5$\times$ the AP of Joint Training (0.018) and 15.8$\times$ that of Naive Sequential (0.004).
The strong performance of CMKL can be attributed to its multimodal feature integration: by leveraging textual (BiomedBERT), molecular (Morgan FP), and structural (R-GCN) node representations through cross-attention fusion, CMKL captures richer entity semantics than embedding-only methods.

\paragraph{EWC reduces forgetting.}
EWC achieves AF\,=\,0.017$\pm$0.001, a 19\% reduction in forgetting compared to Naive Sequential (AF\,=\,0.021$\pm$0.000).
This confirms that Fisher information-based parameter regularization is effective for biomedical KG continual learning, despite the domain's heterogeneity.
The corresponding REM score of 0.984 indicates that EWC retains 98.4\% of previously learned knowledge.

\paragraph{Experience Replay shows limited benefit.}
Surprisingly, Experience Replay achieves AF\,=\,0.021---identical to Naive Sequential.
We hypothesize that the 10\% buffer size is insufficient given the scale of PrimeKG (8.1M+ edges), and that random exemplar selection does not adequately represent the diverse relation types.
CMKL's K-means diverse selection for multimodal memory replay addresses this limitation.

\paragraph{LKGE provides moderate improvement.}
LKGE with TransE achieves AP\,=\,0.039$\pm$0.001 with low forgetting (AF\,=\,0.012$\pm$0.002), demonstrating that its embedding regularization and knowledge distillation are effective.
However, LKGE's reliance on TransE limits its performance on PrimeKG's heterogeneous relation types, where DistMult-based methods perform better.

\paragraph{Joint Training is not the upper bound for AP.}
Counterintuitively, Joint Training achieves only AP\,=\,0.018, lower than both CMKL and LKGE.
This occurs because Joint Training evaluates on the union of all tasks simultaneously, diluting per-task performance, while CMKL and LKGE benefit from task-specific adaptation.
Joint Training does, however, achieve perfect remembering (REM\,=\,1.000) as expected.

\paragraph{RAG agent achieves lowest LP performance.}
The RAG agent (AP\,=\,0.002) is not designed for exhaustive entity ranking and instead excels at targeted QA tasks.
Its near-zero forgetting (AF\,=\,0.001) reflects the non-parametric nature of retrieval-augmented approaches.

\subsection{Node Classification Results}
\label{sec:nc_results}

\Cref{tab:nc_results} presents node classification results for disease categorization.

\begin{table}[t]
\centering
\caption{Node classification results (disease categorization) on \ours (mean $\pm$ std over 5 seeds). AP: Average Performance, AF: Average Forgetting, BWT: Backward Transfer. Best AP is \textbf{bolded}.}
\label{tab:nc_results}
\small
\begin{tabular}{lccc}
\toprule
\textbf{Method} & \textbf{AP} $\uparrow$ & \textbf{AF} $\downarrow$ & \textbf{BWT} $\uparrow$ \\
\midrule
Naive Sequential & $0.344 \pm 0.004$ & $0.011 \pm 0.002$ & $0.003 \pm 0.005$ \\
Joint Training & $0.370 \pm 0.002$ & $0.003 \pm 0.001$ & $0.022 \pm 0.003$ \\
EWC~\citep{Kirkpatrick2017} & $0.345 \pm 0.004$ & $0.008 \pm 0.003$ & $0.007 \pm 0.004$ \\
Experience Replay~\citep{Rolnick2019} & $0.344 \pm 0.006$ & $0.010 \pm 0.003$ & $0.006 \pm 0.005$ \\
\midrule
\textbf{CMKL (ours)} & $\textbf{0.431} \pm 0.005$ & $0.004 \pm 0.002$ & $-0.000 \pm 0.002$ \\
\bottomrule
\end{tabular}
\end{table}


CMKL achieves the highest AP (0.431$\pm$0.005) with minimal forgetting (AF\,=\,0.004$\pm$0.002), demonstrating that multimodal features are particularly beneficial for node classification where entity semantics directly inform the prediction.
Joint Training achieves AP\,=\,0.370$\pm$0.002, and notably all methods show positive BWT for node classification, suggesting that learning new disease associations on later tasks improves classification of previously seen diseases.
EWC, Experience Replay, and Naive Sequential perform similarly (AP $\approx$ 0.344--0.345), indicating that forgetting mitigation provides less benefit for node classification compared to link prediction.

\subsection{KGQA Results}
\label{sec:kgqa_results}

We evaluate the RAG agent (Qwen2.5-7B-Instruct with ChromaDB retrieval) on the link prediction task to provide a non-parametric baseline.
As shown in \cref{tab:main_results}, RAG achieves AP\,=\,0.002$\pm$0.001, the lowest among all methods.
This is expected: the LLM-based approach must generate exact entity names from free-text output, a task for which it is not designed.
However, RAG exhibits near-zero forgetting (AF\,=\,0.001) because retrieval-augmented approaches do not suffer from parametric catastrophic forgetting---the knowledge is stored in the index, not in model weights.
This establishes RAG as a complementary paradigm: while unsuitable for exhaustive entity ranking, it offers stable performance across temporal snapshots and could serve KGQA tasks where natural language answers are acceptable.

\subsection{Analysis}
\label{sec:analysis}

\paragraph{EWC vs.\ Replay forgetting patterns.}
EWC and Experience Replay exhibit different forgetting profiles across entity types.
EWC provides consistent forgetting reduction across all relation types, while Experience Replay's effectiveness varies: it helps most for high-frequency relation types (gene-protein interactions) but provides little benefit for rare types (drug-disease indications), where the 10\% buffer is too small to capture sufficient diversity.

\paragraph{Gene/protein dominance effect.}
Gene/protein entities account for approximately 73\% of triples in PrimeKG, creating a substantial type imbalance.
This affects all methods: performance on gene/protein-centric tasks (Tasks 1--6) is consistently higher than on tasks involving rarer entity types (Tasks 7--10).
The DistMult decoder handles this imbalance better than TransE due to its bilinear scoring function, which can model diverse relation patterns without geometric constraints.

\paragraph{Domain-specific insights.}
The biomedical domain introduces unique CGL challenges absent from generic benchmarks.
First, ontology updates (e.g., MONDO disease hierarchy changes) cause systematic triple removals that look like ``forgetting'' but actually reflect corrected knowledge.
Second, highly connected hub nodes (e.g., TP53, BRCA1) participate in many relation types, creating inter-task dependencies that complicate task-isolated continual learning.
Third, the multimodal nature of biomedical entities means that forgetting can be modality-specific: a model may retain a drug's molecular properties while forgetting its textual description-based associations.
